\guidelinesubsubsection{Recommendations}

When LLMs are used to automate research or software development tasks previously performed by humans, it is essential to assess the LLM's performance. When existing reference datasets or comparison metrics cannot fully capture the target qualities relevant in the study context, researchers \should rely on human judgment to validate the LLM outputs.

Integrating human participants in the study design will require additional considerations, including a recruitment strategy, annotation guidelines, training sessions, or ethical approvals.
Therefore, researchers \should consider human validation early in the study design, not as an afterthought.
Authors \must clearly define the constructs that the human and LLM annotators evaluate~\cite{DBLP:conf/ease/RalphT18}.
When designing custom instruments to assess LLM output (e.g., questionnaires, scales), researchers \should share their instruments in the \supplementarymaterial.

When the judgment is subjective (i.e. depends on the judge's values or implicit or explicit theories of the phenomenon in question), the LLM output \should be evaluated against judgments aggregated from multiple human judges. In these cases, researchers \should clearly describe their aggregation method and document their reasoning. \added{The preferred method is to randomly order the objects requiring judgment and put them in groups small enough to judge in 1--3 hours. Two or three human experts review the first group together, simultaneously judging, discussing, and creating a set of decision rules documenting their reasoning. Then, the experts iterate among rounds of:}
\begin{itemize}
    \item independently rating one group of objects
    \item calculating inter-rater agreement (IRA) or reliability (IRR)
    \item meeting to discuss disagreements and reach consensus
    \item updating the decision rules so the disagreement will not recur.
\end{itemize}

\added{The goal of this process is to reach some target IRA or IRR (e.g. Krippendorff's $\alpha>0.8$), which indicates that the decision rules are sufficient. If the number of judgments required is large, it is permissible to continue with a single rater per judgment once the IRA/IRR target is reached and sustained for 2--3 groups.} 

Researchers \should provide measures of IRA or IRR, preferably broken down by round. Confounding factors \should be controlled for (e.g., by categorizing participants according to their level of experience or expertise).
Where applicable, researchers \should perform a power analysis to estimate the required sample size, ensuring sufficient statistical power in their experimental design.

Researchers \should use established reference models to compare humans with LLMs.
For example, \citeauthor{Schneider2025ReferenceModel}~\cite{Schneider2025ReferenceModel} outline design considerations for studies comparing LLMs with humans.
%They outline the relevant groups (researchers, experimenters, subjects), activities (study planning, task assignment, results analysis, interpretation), and artifacts (study design, tasks, results, conclusion).
%Further, they illustrate the relationship between these entities, the order of steps, and which design decisions need to be considered.
If studies aim to automate annotating software artifacts using an LLM, researchers \should follow systematic approaches to decide whether and how human annotators can be replaced (e.g. showing that a jury of three LLMs exhibit model-to-model agreement exceeding some threshold).  
% For example, \citeauthor{DBLP:conf/msr/AhmedDTP25}~\cite{DBLP:conf/msr/AhmedDTP25} suggest a method that involves using a jury of three LLMs with 3 to 4 few-shot examples rated by humans, where the model-to-model agreement on all samples is determined using \emph{Krippendorff's Alpha}.
% If the agreement is high (alpha > 0.5), a human rating can be replaced with an LLM-generated one.
% In cases of low model-to-model agreement (alpha $\le$ 0.5), they then evaluate the prediction confidence of the model, selectively replacing annotations where the model confidence is high ($\ge$ 0.8).

Besides generating and modifying content, agentic software development tools such as Claude Code can autonomously call command-line tools or pull in additional information from MCP servers.
For such tools, it makes sense to separate textual output (e.g., file changes) from output related to tool execution.
When evaluating agentic tools, researchers \should assess the feedback that users provided for proposed changes, report statistics on how frequently they accepted the content, and how they modified it.
This agentic human-in-the-loop interaction approach is essentially a built-in human validation, even though the degrees of freedom are larger than in more traditional experiments that use LLMs directly.

\guidelinesubsubsection{Example(s)}

\citeauthor{DBLP:conf/msr/AhmedDTP25}~\cite{DBLP:conf/msr/AhmedDTP25} evaluated how human annotations can be replaced by LLM generated labels.
In their study, they explicitly report the calculated agreement metrics between the models and between humans. They found that three LLMs (GPT-4, Gemini, and Claude) had better agreement on some tasks than human annotators. Meanwhile, \citeauthor{DBLP:conf/icse/XueCBTH24}~\cite{DBLP:conf/icse/XueCBTH24} conducted a controlled experiment to evaluate the impact of ChatGPT on the performance and perceptions of students in an introductory programming course.
They used multiple measures to judge LLM impacts from the human point of view including recording students' screens, evaluating their answers for given tasks, and distributing an opinion survey.

%\citeauthor{hymel2025analysisllmsvshuman}~\cite{hymel2025analysisllmsvshuman} evaluated the capability of ChatGPT-4.0 to generate requirements documents. 
%Specifically, they evaluated two requirements documents based on the same business use case, one document generated with the LLM and one document created by a human expert.
%The documents were then reviewed by experts and judged in terms of alignment with the original business use case (scale of 1-10), requirements completion (Not Complete, Fairly Complete, Fully Complete), and whether they believed it was created by a human or an LLM.
%Finally, they analyzed the influence of the participants' familiarity with AI tools on the study results.
%Participants self-reported their AI tool proficiency (Novice, Intermediate, Advanced, Expert) and usage frequency (Daily, Weekly, Monthly, Never), which were then correlated with their ratings of the requirements documents.

\guidelinesubsubsection{Benefits}

Validating LLMs against human judgments builds confidence in the LLM's accuracy and validity, increasing the trustworthiness of findings, especially when IRA or IRR metrics are reported~\cite{khraisha2024canlargelanguagemodelshumans}. Furthermore, incorporating feedback from human judges may help to improve the LLM or LLM-based tool based on the reported experiences.


%For example, LLM results vary significantly in natural language processing tasks, casting doubt on their reliability and the possibiltiy of replacing human judges~\cite{DBLP:journals/corr/abs-2406-18403}. Moreover, LLMs do not match human annotation in labeling tasks for natural language inference, position detection, semantic change, and hate speech detection~\cite{DBLP:conf/chi/Wang0RMM24}.



\guidelinesubsubsection{Challenges}

Assessing a construct like LLM performance is challenging because ensuring that a construct is defined well and operationalized using an appropriate measurement model requires a deep understanding of (1) the construct, (2) construct validity in general, and (3) instrumentation~\cite{DBLP:journals/tse/SjobergB23,ralph2024teaching}. Comparing an LLM to human judges is typically slower and more expensive than machine-generated measures. More fundamentally, neither human judgment nor machine-generated measures provides an objective ground truth against which LLM accuracy can be firmly determined.   

Human judgments exhibit variability due to differences in experience, expertise, interpretations, and personal biases~\cite{DBLP:journals/pacmhci/McDonaldSF19}. When diverse humans---given clear decision rules corresponding to our construct definitions---rate items reliably, \textit{we simply assume that their reliability implies measurement validity}. It does not. We assume that reliability implies validity because measuring reliability is \textit{much} easier than measuring validity. Indeed, much of the time the best researchers can do is make a conceptual argument for why these judges armed with these decision rules based on this construct \textit{should} produce valid ratings. 


%For example, \citeauthor{hicks_lee_foster-marks_2025} found that racialized software developers \enquote{rated the quality of the output of AI-assisted coding tools significantly lower than other groups}~\cite{hicks_lee_foster-marks_2025}.
% Such biases can be mitigated with a careful selection and combination of multiple judges, but cannot be completely controlled for.

Researchers must weigh the cost and time need for human validation against the expected corresponding improvement in validity over machine-generated measures. 

\guidelinesubsubsection{Study Types}

For \llmusage, researchers \should carefully reflect on the validity of their evaluation criteria.
When an extensively-validated and widely-used benchmark is available (see \benchmarkingtasks), there is little need for human validation.
Where criteria are less clear (e.g., \annotators, \subjects, or \synthesis), or researchers are creating and validating a new benchmark, human validation is important, and this guideline should be followed.
When using \judges, it is common to start with humans who co-create initial rating criteria. 
In developing \newtools, human validation of the LLM output is critical when the output is supposed to match human expectations.

\guidelinesubsubsection{Advice for Reviewers}

\added{Human validation may be the most difficult of our guidelines for reviewers to assess because it will often involve cutting through fallacious arguments. Authors have strong incentives to avoid human validation. However, if LLM output is validated only by comparison with the same or other LLMs, reviewers should demand \textit{quantitative empirical evidence} that such comparison is reliable \textit{and} valid. Similarly, reviewers should demand \textit{quantitative empirical evidence} that any employed benchmarks are reliable \textit{and} valid. Without such evidence, insist on human validation. Showing that several LLMs exhibit high agreement is not sufficient because reliability is necessary but insufficient for validity.}

\added{Next, reviewers will encounter dubious approaches to human validation, such as having a single judge (typically an author) or recruiting dozens of anonymous, untrained, inexperienced gig-workers and aggregating their results on a majority-vote basis. Reviewers should only accept a single judge when the judgments depend only on widely-accepted theories and entails limited value conflict (e.g. tagging method names that contain abbreviations or acronyms). Unfortunately, researchers tend to underestimate the level how theory- and value-laden phenomena are. Any assessment of code quality or tagging of bugs, for example, would be highly dependent on the judge's values and implicit theories~\cite{ralph2024teaching}. }

\added{For multiple judges, reviewers should expect research designs that include techniques known to improve IRA/IRR such as recruiting raters with appropriate education and experience, organizing the rating task into rounds, having raters conduct the initial round of judgments together, having consensus meetings (not voting) after each round in which decision rules are updated. Reviewers should not accept low IRR or IRA (e.g. Krippendorff's $\alpha<0.8$) for research designs absent any of these techniques. Conversely, if authors have embraced best practices for improving IRR/IRA and still obtained middling results (e.g. $0.66<\alpha<0.8$) reviewers should merely insist that low reliability be noted as a limitation.} 

\added{Moving beyond reliability, reviewers should expect authors to \textit{explain} (conceptually, not qualitatively) why human judgments should be valid by comparing construct definitions, directions, decision rules, and judge expertise.}

\added{As with previous guidelines, when important information is missing, reviewers should simply request it be added unless so much is missing that the reviewer cannot assess methodological rigor. For example, failing to provide judge instructions, instruments, decision rules, or construct definitions may prevent assessment of rigor and validity, while missing information about recruitment is less serious. Reviewers should be patient when requesting clearer definitions---the same construct definitions that seem clear to authors often seem opaque to reviewers, and papers should not be rejected over routine clarification requests.}